% =============================================================
%  chapter5_data_table.tex -- AUTO-GENERATED, DO NOT EDIT BY HAND.
%  Written by Chapter5/Figures/write_chapter5_data_table.py from
%  Chapter5/Tables/Draaikraal 2023 BJT.xlsx ('Geotagged Stratigraphy').
%  Regenerate after any change to the spreadsheet:
%    python3 Chapter5/Figures/write_chapter5_data_table.py
%  \input{} by Chapter5/chapter5.tex (defines the macro); invoked from
%  Chapter5/Writing/chapter5_SI_body.tex.
% =============================================================
\providecommand{\ChFivePhoto}[1]{#1}
\newcommand{\ChFiveFieldStationTable}{%
\begingroup
\edef\ChFiveSavedFootnote{\the\value{footnote}}\setcounter{footnote}{0}
\renewcommand{\thefootnote}{\arabic{footnote}}
\begin{scriptsize}
\fontsize{7.5}{8.6}\selectfont
\renewcommand{\arraystretch}{0.9}
\setlength{\tabcolsep}{3pt}
\begin{longtable}{@{}>{\raggedright\arraybackslash}p{1.95cm} r r r r l l l >{\raggedright\arraybackslash}p{1.5cm}@{}}
\caption{Field stations of the Draaikraal campaign (July 2023), from the `Geotagged Stratigraphy' sheet of \texttt{Chapter5/Tables/Draaikraal 2023 BJT.xlsx}, in order of stratigraphic height. Ht.: stratigraphic height (m) relative to the Merensky Reef, from the projection described in \S\ref{sec:ch5-stratigraphic-column-methods}. Lat./Lon.: WGS84, decimal degrees. Elev.: elevation (m). Unit: final interpreted map unit and Lith.: lithology code (\S\ref{sec:ch5-lithologic-descriptions}). Structure: bedding (Bed.), cleavage (Clv.), contact (Cont.) or modal layering (M.L.), as strike/dip with strike an azimuth from north. Photos: numbers in the photo album (\S\ref{sec:ch5-SI-photo-album}), each a link. Field descriptions and notes are footnotes to the station name, at the foot of the page. `--': not recorded.}\label{tab:ch5-field-stations}\\
\toprule
Station & Ht. & Lat. & Lon. & Elev. & Unit & Lith. & Structure & Photos \\
\midrule
\endfirsthead
\multicolumn{9}{@{}l}{\textit{Table~\ref{tab:ch5-field-stations} (continued)}} \\
\toprule
Station & Ht. & Lat. & Lon. & Elev. & Unit & Lith. & Structure & Photos \\
\midrule
\endhead
\midrule
\multicolumn{9}{r@{}}{\textit{continued on next page}} \\
\endfoot
\bottomrule
\endlastfoot
RZ22 & -933 & -25.24760 & 30.11812 & 2146 & Rz & Msp & -- & -- \\
QZ27 & -724 & -25.24040 & 30.10576 & 2076 & Rz & Qz & -- & -- \\
QZ7 & -706 & -25.23020 & 30.11221 & 1958 & Rz & Qz & -- & -- \\
RZ54 & -700 & -25.23980 & 30.10496 & 2044 & Rz & Msp & -- & -- \\
CONT MARG E Q W?\footnote{\scriptsize Quartzite. Quartzite follows topo line locally here.} & -697 & -25.22880 & 30.11288 & 1960 & Rz & Qz & -- & -- \\
D3NB & -683 & -25.22540 & 30.11481 & 1960 & Rz & Qz & Bed. 145/30 & -- \\
QZ1 & -676 & -25.22530 & 30.11455 & 1956 & Rz & Qz & -- & -- \\
D3NA\footnote{\scriptsize Compass trouble, resolved.} & -674 & -25.22590 & 30.11390 & 1929 & Rz & Qz & Bed. 120/30 & -- \\
RZ43 & -669 & -25.23920 & 30.10392 & 2015 & Rz & Msp & -- & -- \\
D7NC\footnote{\scriptsize Alternating quartzite-marginal zone photo, not very good one.} & -661 & -25.23900 & 30.10369 & 2010 & Rz & Qz & -- & -- \\
RZ8 & -654 & -25.22500 & 30.11369 & 1945 & Rz & Msp & -- & -- \\
D3NC\footnote{\scriptsize No faulting; the quartz and marginal zone is interbedded.} & -649 & -25.22510 & 30.11337 & 1940 & Rz & QzMsp & -- & \ChFivePhoto{34} \\
RZ21 & -644 & -25.22500 & 30.11326 & 1947 & Rz & Msp & -- & -- \\
RZ24 & -643 & -25.23900 & 30.10279 & 1997 & Rz & Msp & -- & -- \\
RZ36 & -642 & -25.23860 & 30.10304 & 1995 & Rz & Msp & -- & -- \\
QZ112 & -639 & -25.23860 & 30.10289 & 1995 & Rz & Qz & -- & -- \\
D3SF\footnote{\scriptsize Qz-Px Pegmatite. Hosted in migmatitmatic quartzite. Schist banding deformed around pegmatite.} & -638 & -25.22610 & 30.11207 & 1929 & Rz & ScP & -- & \ChFivePhoto{37}, \ChFivePhoto{38} \\
SCHIS & -636 & -25.22580 & 30.11224 & 1937 & Rz & Sc & -- & -- \\
QZ2 & -631 & -25.22510 & 30.11260 & 1950 & Rz & Qz & -- & -- \\
QZ3 & -631 & -25.22460 & 30.11297 & 1954 & Rz & Qz & -- & -- \\
RZ11 & -629 & -25.22490 & 30.11266 & 1952 & Rz & Msp & -- & -- \\
QZ6 & -629 & -25.22520 & 30.11238 & 1946 & Rz & Qz & -- & -- \\
QZ5 & -626 & -25.22480 & 30.11255 & 1950 & Rz & Qz & -- & -- \\
RZ31 & -625 & -25.22600 & 30.11153 & 1926 & Rz & Msp & -- & -- \\
RZ33 & -622 & -25.22460 & 30.11252 & 1946 & Rz & Msp & -- & -- \\
D3SE\footnote{\scriptsize Fine grained, light rock. Dacite? Silicified. Pockmarks weathering proud in paleovertical pipes.} & -621 & -25.22460 & 30.11251 & 1950 & Rz & Msp & -- & \ChFivePhoto{36}, \ChFivePhoto{40} \\
D3SD\footnote{\scriptsize Migmatized quartzite. Banded. Light bands of quartz preserve original sedimentary texture. Grains rounded and 2 mm. Some bands pink (feldspar) spme green (chlorite, epidote?)} & -619 & -25.22450 & 30.11247 & 1946 & Rz & ScQz & -- & \ChFivePhoto{35} \\
RZ71 & -617 & -25.22690 & 30.11044 & 1911 & Rz & Msp & -- & -- \\
SCHIST & -615 & -25.22440 & 30.11236 & 1944 & Rz & Sc & -- & -- \\
QZ4 & -615 & -25.22430 & 30.11241 & 1945 & Rz & Qz & -- & -- \\
RZ41 & -614 & -25.22410 & 30.11254 & 1948 & Rz & Msp & -- & -- \\
RZ110 & -614 & -25.23840 & 30.10183 & 1976 & Rz & Msp & -- & -- \\
RZ51 & -614 & -25.22290 & 30.11342 & 1945 & Rz & Msp & -- & -- \\
RZ61 & -599 & -25.22550 & 30.11068 & 1907 & Rz & Msp & -- & -- \\
RZ32 & -588 & -25.22590 & 30.10982 & 1900 & Rz & Msp & -- & -- \\
QZ26 & -575 & -25.23760 & 30.10054 & 1948 & Rz & Qz & -- & -- \\
D7SE\footnote{\scriptsize Silicate Pegmatite. within migmatized qzite very close to magnesite skarn. Mymerkitic quartz-K-spar, huge OPX, small vibrant rare green and blue crystals: apatite, CPX, beryl?} & -569 & -25.23690 & 30.10081 & 1946 & Rz & ScP & -- & \ChFivePhoto{70}, \ChFivePhoto{71} \\
SHIST & -569 & -25.23690 & 30.10078 & 1947 & Rz & ScP & -- & -- \\
RZUP & -565 & -25.23740 & 30.10019 & 1935 & Rz & ScPMsp & -- & -- \\
D7SD\footnote{\scriptsize Magnesite Skarn. within migmatized qzite. In ground, hard to tell if dike or concordant with schist host.} & -559 & -25.23740 & 30.09991 & 1936 & Rz & ScP & -- & \ChFivePhoto{69} \\
D2NC\footnote{\scriptsize Contact. On this ridge, there are multiple bands of quartzites with Marginal Zone between them. Is the marginal zone a sill between layers of quartzite, or is there a fault? (Update, I think it’s a strike-slip(?) fault striking E-W dipping to the South with a component of left-lateral shear along the ridge)} & -555 & -25.22110 & 30.11222 & 1993 & Rz & QzMsp & -- & \ChFivePhoto{14}, \ChFivePhoto{23}, \ChFivePhoto{28} \\
D7NA\footnote{\scriptsize Fine marginal gabbronorite with shearing pegmatite veins. sheared pegmatite veins in pockets/pipes orthogonal to bedding, paleoup. This layer is just above migmatized quartzite and just below fine magnetic marginal zone. Represents fluid transport from migmatized sedimentary rocks, origin of resistance banding? Strike and Dip are for migmatized Qz.} & -552 & -25.23700 & 30.09985 & 1924 & Rz & QzMsp & Bed. 220/25 & \ChFivePhoto{67} \\
SCH1 & -551 & -25.23700 & 30.09984 & 1930 & Rz & ScP & -- & -- \\
RZMAG & -549 & -25.23690 & 30.09979 & 1919 & Rz & Msp & -- & -- \\
RZ & -532 & -25.22390 & 30.10878 & 1909 & Rz & Msp & -- & -- \\
D1ND\footnote{\scriptsize Bladed iridescent white crystals (Plagioclase) in 2 cm thick planar features intersect orthogonally in outcrop. The plagioclase planes are surrounded by 10 cm zone of oikocrystic augite. The southwestern limb pinches in and out and occasionally overlaps itself in parallel, a la boudinage.} & -507 & -25.23730 & 30.09753 & 1926 & Rz & Msp & -- & \ChFivePhoto{6} \\
D1NC\footnote{\scriptsize Note that all marginal zone rocks are characterized by dark resistance bands, perhaps cataclastic, reflecting warm-brittle deformation with finely recrystallized zones. They stand proud. They frequently strike E-W, secondarily, NE-SW, tertiarily.  NW-SE, but can vary and appear pretty random. This feature is diagnostic of the marginal zone.} & -500 & -25.23720 & 30.09730 & 1924 & Rz & Msp & -- & -- \\
RZ3 & -485 & -25.22250 & 30.10765 & 1910 & Rz & Msp & -- & -- \\
RLS & -481 & -25.22350 & 30.10665 & 1890 & Cz & Pg & -- & -- \\
RZ1 & -481 & -25.22030 & 30.10931 & 1973 & Rz & Msp & -- & -- \\
D1Sb\footnote{\scriptsize Common Marginal zone Medium-fine grained gabbronorite.} & -479 & -25.23690 & 30.09644 & 1891 & Rz & Msp & -- & \ChFivePhoto{19} \\
D1A\footnote{\scriptsize Fine-grained Magnetite Gabbro with gabbronoritic blocks. Fine-grained magnetite gabbro includes rounded, elongate blocks of oikocrystic gabbronorite preferentially horizontally imbricated, $\sim$6-10 cm along axis. Also a bunch of rounded gabbronorite inclusions that may represent the spherical aggregates of the TBM.} & -479 & -25.23690 & 30.09644 & 1891 & Rz & MspTB & -- & \ChFivePhoto{1}, \ChFivePhoto{4} \\
MZ3 & -478 & -25.22290 & 30.10699 & 1902 & Cz & Pg & -- & -- \\
D2SB\footnote{\scriptsize Sampled because quartzite here all have iron stained interiors. The weathered exteriors are clear. Quartzites served as a conduit for mobilized iron-rich fluids?} & -477 & -25.21950 & 30.10983 & 2004 & Rz & Qz & -- & -- \\
RZ2 & -470 & -25.22270 & 30.10679 & 1901 & Rz & Msp & -- & -- \\
D2SH\footnote{\scriptsize Schist with Skarn. Swirled and deformed schist with marginal zone, with thin clear banding. Non-magnetic. Xenolith? Associated with skarn pegmatite veins: 60\% Feldspar, 40\% amphibole? Big very dark crystals that I think I see 120 degree cleavage on… Likely devolatilizing xenolith.} & -469 & -25.22070 & 30.10839 & 1950 & Rz & ScP & -- & \ChFivePhoto{29}, \ChFivePhoto{30} \\
RQCONT & -465 & -25.21920 & 30.10955 & 2014 & Rz & QzMsp & -- & -- \\
RZ20 & -456 & -25.23540 & 30.09654 & 1903 & Rz & Msp & -- & -- \\
RZ5 & -441 & -25.23760 & 30.09423 & 1918 & Rz & Msp & -- & -- \\
D2ND\footnote{\scriptsize Measureable bedding planes at top of ridge: Strike: 265, Dip: 25. Also, One bedding plane preserved lovely sedimentary ripples striking E-W.} & -440 & -25.21820 & 30.10918 & 2020 & Rz & Qz & Bed. 175/25 & \ChFivePhoto{20} \\
D11SA\footnote{\scriptsize Coarse grained gabbronorite. Is it critical-margin zone? Are resistance bands chromite bearing?} & -440 & -25.23770 & 30.09405 & 1899 & Rz & Msp & -- & -- \\
QZ FOOLLOW TOPO & -433 & -25.21900 & 30.10814 & 1990 & Rz & Qz & -- & -- \\
RZ4 & -431 & -25.23700 & 30.09419 & 1906 & Rz & Msp & -- & -- \\
RZ6 & -430 & -25.23860 & 30.09298 & 1923 & Rz & Msp & -- & -- \\
D2NE\footnote{\scriptsize Different strike, same dip. Strike: 315, dip: 25. Quartzites are rotated around the top of the ridge. NW-SE vertical fault or E-W South dipping fault can neatly explain rotation.} & -426 & -25.21670 & 30.10968 & 2031 & Rz & Qz & Bed. 225/25 & -- \\
D4NB\footnote{\scriptsize Contact between floor quartzite and marginal zone. Bed plane of quartzite measured.} & -422 & -25.24150 & 30.09048 & 1932 & Rz & QzMsp & Bed. 240/25 & \ChFivePhoto{52} \\
QZ51 & -417 & -25.21880 & 30.10754 & 1993 & Rz & Qz & -- & -- \\
D10NK & -415 & -25.24210 & 30.08966 & 1917 & Rz & Qz & Bed. 235/25 & -- \\
RZ7 & -413 & -25.23870 & 30.09213 & 1915 & Rz & Msp & -- & -- \\
D4SA\footnote{\scriptsize Contact between marginal zone and migmatized quartzite. Migmatitic quartzite $\sim$10m thick. At contact is non-magnetic leucocratic gneiss, then fine magnetic marginal zone: 10cm, then coarse grained marginal zone: 30 cm, then common fine marginal zone.} & -412 & -25.24090 & 30.09036 & 1895 & Rz & QzMsp & -- & \ChFivePhoto{50}, \ChFivePhoto{51} \\
SCHI & -403 & -25.24040 & 30.09025 & 1882 & Rz & Sc & -- & -- \\
QZ61 & -402 & -25.21890 & 30.10676 & 1980 & Rz & Qz & -- & -- \\
QZ41 & -402 & -25.21870 & 30.10691 & 1987 & Rz & Qz & -- & -- \\
QZ & -395 & -25.21680 & 30.10808 & 2001 & Rz & Qz & -- & -- \\
QZ29 & -393 & -25.24210 & 30.08861 & 1907 & Rz & Qz & -- & -- \\
QZ71 & -391 & -25.21900 & 30.10613 & 1975 & Rz & Qz & -- & -- \\
RZ28 & -388 & -25.24200 & 30.08845 & 1906 & Rz & Msp & -- & -- \\
SCH & -387 & -25.24010 & 30.08972 & 1870 & Rz & Sc & -- & -- \\
RZ12 & -381 & -25.23990 & 30.08955 & 1866 & Rz & Msp & -- & -- \\
D2SA\footnote{\scriptsize Common Main Zone Coarse Grained Gabbronorite (Office interp: Leucogabbronorite; porphyritic gabbronorite) On Dirt road to Klipriver Farms, Boulder piles outcropping of main zone. These rocks are the building materials for much of the structures in the area. Weathering is quite distinct from Marginal zone. Both have exfoliation/onion skin weathering. Main Zone is prouder, forming impressive rounded boulders like a granite. Weathered faces have distinct dalmatian pattern. A large boulder (float) at this outcrop contains a characteristic pyroxene “tennisball”, a well-known marker at the base of the Main Zone.} & -380 & -25.23000 & 30.09696 & 1860 & Cz & PgTB & -- & -- \\
RZ35 & -380 & -25.21890 & 30.10566 & 1965 & Rz & Msp & -- & -- \\
D2SF\footnote{\scriptsize Coarse Grained Magnetite Gabbro. Actually I am now thinking this sample is similar to the magnetite gabbro D1Sa, but coarser grained. That sample too had zones of coarse grains, and this one also had blocks of oikocrystic augite.} & -379 & -25.21550 & 30.10829 & 1994 & Rz & Msp & -- & \ChFivePhoto{24} \\
RZ42 & -379 & -25.21940 & 30.10519 & 1950 & Rz & Msp & -- & -- \\
QZ8 & -377 & -25.21920 & 30.10529 & 1954 & Rz & Qz & -- & -- \\
RZ9 & -370 & -25.23920 & 30.08952 & 1850 & Rz & Msp & -- & -- \\
RAT & -368 & -25.23630 & 30.09174 & 1883 & Rz & Msp & -- & -- \\
RZ511 & -368 & -25.21970 & 30.10442 & 1935 & Rz & Msp & -- & -- \\
RZLQZR & -366 & -25.21960 & 30.10440 & 1938 & Rz & QzMsp & -- & -- \\
RZ53 & -362 & -25.21930 & 30.10448 & 1945 & Rz & Msp & -- & -- \\
RZ113 & -359 & -25.24170 & 30.08722 & 1876 & Rz & Msp & -- & -- \\
IS1 & -358 & -25.23030 & 30.09577 & 1868 & IS & IS & -- & -- \\
RZ62 & -355 & -25.21960 & 30.10390 & 1934 & Rz & Msp & -- & -- \\
D2SG\footnote{\scriptsize Norite. Weathering more closely resembles main zone than nearby marginal zone, but this rock does not look like main zone to me. Non-magnetic.} & -352 & -25.21390 & 30.10811 & 1956 & Cz & Dcg & -- & \ChFivePhoto{27} \\
QZ9 & -352 & -25.22000 & 30.10342 & 1929 & Rz & Qz & -- & -- \\
SCI & -350 & -25.23900 & 30.08890 & 1888 & Rz & Sc & -- & -- \\
D9NE & -348 & -25.25130 & 30.07965 & 1929 & Rz & Qz & Bed. 180/15 & -- \\
QZ10 & -347 & -25.22000 & 30.10320 & 1928 & Rz & Qz & -- & -- \\
SCH11 & -343 & -25.24070 & 30.08725 & 1881 & Rz & Sc & -- & -- \\
RZ27 & -341 & -25.24060 & 30.08725 & 1882 & Rz & Msp & -- & -- \\
D9SD\footnote{\scriptsize Fine Grained Marginal Norite. Sampled for prominent resistance band for later analyses.} & -340 & -25.25150 & 30.07907 & 1912 & Rz & Msp & -- & \ChFivePhoto{95} \\
QZ11 & -339 & -25.21970 & 30.10308 & 1938 & Rz & Qz & -- & -- \\
SCH2 & -338 & -25.23940 & 30.08802 & 1887 & Rz & Sc & -- & -- \\
RZ82 & -338 & -25.23960 & 30.08784 & 1886 & Rz & Msp & -- & -- \\
D9NC\footnote{\scriptsize Chloritized Norite. Plag, OPX, abundant phlogopite, alteration grade lessens up section.} & -335 & -25.25180 & 30.07858 & 1908 & Rz & Msp & -- & -- \\
QZ12 & -334 & -25.21970 & 30.10285 & 1939 & Rz & Qz & -- & -- \\
D9NF\footnote{\scriptsize Fine Grained Marginal norite. Droplet forms of coarse pyroxenite xenoliths, pegmatite inclusions?} & -333 & -25.25050 & 30.07951 & 1918 & Rz & MspTB & -- & \ChFivePhoto{96} \\
RZ23 & -331 & -25.21720 & 30.10466 & 1957 & Rz & Msp & -- & -- \\
RZ72 & -330 & -25.21940 & 30.10291 & 1940 & Rz & Msp & -- & -- \\
RZ81 & -318 & -25.21920 & 30.10248 & 1940 & Rz & Msp & -- & -- \\
RZ19 & -316 & -25.21700 & 30.10410 & 1948 & Rz & Msp & -- & -- \\
RZ25 & -312 & -25.24970 & 30.07913 & 1910 & Rz & Msp & -- & -- \\
D9NG\footnote{\scriptsize Paleovertical Pegmatite Dike Photograph, Marginal zone coarsens upsection towards quartzite.} & -312 & -25.24980 & 30.07904 & 1910 & Rz & Msp & -- & \ChFivePhoto{97}, \ChFivePhoto{98} \\
D5NC\footnote{\scriptsize Cleavage.} & -299 & -25.21880 & 30.10195 & 1946 & Rz & Msp & Clv. 145/25 & \ChFivePhoto{57} \\
RZ44 & -290 & -25.24030 & 30.08497 & 1844 & Rz & Msp & -- & -- \\
RZ18 & -286 & -25.21730 & 30.10245 & 1940 & Rz & Msp & -- & -- \\
RZ91 & -280 & -25.21900 & 30.10086 & 1928 & Rz & Msp & -- & -- \\
QZ31 & -264 & -25.21640 & 30.10209 & 1933 & Rz & Qz & -- & -- \\
QZ28 & -252 & -25.24750 & 30.07796 & 1893 & Rz & Qz & -- & -- \\
RZ10 & -252 & -25.21880 & 30.09972 & 1929 & Rz & Msp & -- & -- \\
QZ13 & -248 & -25.21950 & 30.09891 & 1907 & Rz & Qz & -- & -- \\
D5SB\footnote{\scriptsize Amphibole Gabbro. Weakly magnetic, amphibole bearing? Rock clearly rotten.} & -248 & -25.21430 & 30.10287 & 1919 & Cz & A & -- & -- \\
D9NH & -244 & -25.24740 & 30.07764 & 1895 & Rz & Qz & Bed. 230/35 & -- \\
QZ23 & -242 & -25.21610 & 30.10128 & 1932 & Rz & Qz & -- & -- \\
QZ14 & -241 & -25.21960 & 30.09852 & 1907 & Rz & Qz & -- & -- \\
QZ15 & -233 & -25.21970 & 30.09807 & 1907 & Rz & Qz & -- & -- \\
RPEG1 & -229 & -25.21390 & 30.10229 & 1907 & Cz & Dcg & -- & -- \\
RZ111 & -223 & -25.21950 & 30.09777 & 1902 & Rz & Msp & -- & -- \\
D7NF\footnote{\scriptsize Pegmatitic gabbronorite. large (5 cm) oikocrystic pyroxene.} & -218 & -25.21350 & 30.10206 & 1906 & Cz & Dcg & Bed. 180/10 & \ChFivePhoto{73} \\
QZ110 & -215 & -25.21610 & 30.10002 & 1923 & Rz & Qz & -- & -- \\
QZ113 & -214 & -25.24740 & 30.07621 & 1874 & Rz & Qz & -- & -- \\
RZ121 & -213 & -25.21890 & 30.09779 & 1911 & Rz & Msp & -- & -- \\
QZ16 & -213 & -25.21910 & 30.09763 & 1908 & Rz & Qz & -- & -- \\
QZ17 & -211 & -25.21880 & 30.09778 & 1915 & Rz & Qz & -- & -- \\
RZ13 & -209 & -25.21860 & 30.09780 & 1909 & Rz & Msp & -- & -- \\
QZ18 & -205 & -25.21850 & 30.09772 & 1914 & Rz & Qz & -- & -- \\
D9NI\footnote{\scriptsize Schist (Pebble marker?) well mixed clasts of schists, zones of pink felsic melt.} & -204 & -25.24670 & 30.07624 & 1865 & Rz & ScMsp & -- & \ChFivePhoto{99}, \ChFivePhoto{100}, \ChFivePhoto{101}, \ChFivePhoto{102} \\
QZ19 & -200 & -25.21860 & 30.09741 & 1912 & Rz & Qz & -- & -- \\
RZ112 & -199 & -25.24650 & 30.07613 & 1865 & Rz & Msp & -- & -- \\
D7NG\footnote{\scriptsize Pegmatitic gabbronorite. Here magnetic, slightly upsection ($\sim$2 m) rocks non-magnetic. Large (30 cm) blobs of pyroxenite rimmed by anorthosite. Pyroxene/Anorthosite ranging from 1/4 to 4. Included blocks with preserve reistance banding.} & -197 & -25.21220 & 30.10208 & 1905 & Cz & Dcg & Clv. 230/15 & \ChFivePhoto{72}, \ChFivePhoto{73}, \ChFivePhoto{74}, \ChFivePhoto{75} \\
SCHI1 & -193 & -25.24620 & 30.07609 & 1863 & Rz & Sc & -- & -- \\
RZQZLEFT & -193 & -25.21920 & 30.09654 & 1887 & Rz & QzMsp & -- & -- \\
QZ20 & -187 & -25.21880 & 30.09663 & 1901 & Rz & Qz & -- & -- \\
RZ14 & -185 & -25.21870 & 30.09657 & 1900 & Rz & Msp & -- & -- \\
RZ15 & -175 & -25.21840 & 30.09636 & 1895 & Rz & Msp & -- & -- \\
RZ26 & -173 & -25.24550 & 30.07565 & 1856 & Rz & Msp & -- & -- \\
MZ11 & -166 & -25.22800 & 30.08851 & 1837 & Cz & Pg & -- & -- \\
D5SA\footnote{\scriptsize Thought it looked like the "dark facies" main zone in the field,. Compare to D25g (porhpyitic gabbronorite), Quartzite in float (alluvial)} & -166 & -25.21540 & 30.09820 & 1903 & Cz & UD & -- & -- \\
D6SL\footnote{\scriptsize Porphyritic gabbronorite. In notes, described as main zone porphyritic gabbronorite. However these rocks are struturally beneath chrome layers. Closer inspection reveals that pyroxenes are often oikocrystic and large... Incompletely formed Mottled Anorthosite?} & -165 & -25.21490 & 30.09851 & 1892 & Cz & Pg & -- & -- \\
RZ16 & -162 & -25.21870 & 30.09548 & 1887 & Rz & Msp & -- & -- \\
QZ22 & -161 & -25.21810 & 30.09591 & 1898 & Rz & Qz & -- & -- \\
D5NE & -160 & -25.21870 & 30.09537 & 1883 & Rz & Msp & Clv. 125/25 & -- \\
RZ17 & -154 & -25.21950 & 30.09445 & 1875 & Rz & Msp & -- & -- \\
QZ21 & -153 & -25.21950 & 30.09444 & 1875 & Rz & Qz & -- & -- \\
QZ25 & -147 & -25.21770 & 30.09553 & 1877 & Rz & Qz & -- & -- \\
D6SG\footnote{\scriptsize Magnetite pegmatite. Host rock magnetic. Pegmatite weathers proud, cutting the rock paleo-up, "sandy" texture in host rock.} & -146 & -25.21840 & 30.09491 & 1872 & Rz & QzMsp & -- & \ChFivePhoto{60} \\
D6NK\footnote{\scriptsize Porphyritic gabbronorite.} & -145 & -25.21500 & 30.09749 & 1888 & Cz & Pg & Clv. 120/5 & -- \\
D7SH\footnote{\scriptsize Sulfide-pyroxenite.} & -144 & -25.21020 & 30.10109 & 1898 & Cz & Px & -- & \ChFivePhoto{78} \\
D6SH\footnote{\scriptsize Coarse-grained Feldspatic Pyroxenite. Non-magnetic, Floor of critical zone? "Flatreef"?} & -144 & -25.21770 & 30.09534 & 1871 & Cz & Dcg & -- & \ChFivePhoto{62} \\
D6SI\footnote{\scriptsize Coarse-grained Feldspatic Pyroxenite. Coarser grained rocks here are more rotten.} & -144 & -25.21770 & 30.09534 & 1871 & Cz & Dcg & -- & \ChFivePhoto{61} \\
QZ111 & -140 & -25.21760 & 30.09526 & 1869 & Rz & Qz & -- & -- \\
QZ24 & -139 & -25.21800 & 30.09487 & 1869 & Rz & Qz & -- & -- \\
RPEQ & -135 & -25.21480 & 30.09718 & 1883 & Cz & Dcg & -- & -- \\
PEG & -133 & -25.21140 & 30.09968 & 1893 & Cz & Dcg & -- & -- \\
D6NJ\footnote{\scriptsize Pegmatite gabbronorite.} & -132 & -25.21540 & 30.09655 & 1880 & Cz & Dcg & Clv. 150/10 & \ChFivePhoto{63} \\
RPEG & -130 & -25.21540 & 30.09644 & 1877 & Cz & Dcg & -- & -- \\
RZ34 & -124 & -25.21910 & 30.09334 & 1858 & Rz & Msp & -- & -- \\
RZ37 & -120 & -25.24380 & 30.07447 & 1840 & Rz & Msp & -- & -- \\
RZ52 & -118 & -25.22040 & 30.09203 & 1846 & Rz & Msp & -- & -- \\
D7NI\footnote{\scriptsize Pyroxene tennisballs in Gabbronorite. Structurally below chromite?} & -100 & -25.20893 & 30.09996 & 1876 & Cz & PgTB & -- & \ChFivePhoto{79}, \ChFivePhoto{119} \\
D7SJ\footnote{\scriptsize Gabbronorite. small amount of interstitial CPX. Transitional zone?} & -100 & -25.20890 & 30.09995 & 1876 & Cz & Pg & -- & -- \\
MZPOR & -96 & -25.20910 & 30.09963 & 1873 & Cz & Pg & -- & -- \\
CZ3 & -94 & -25.20880 & 30.09978 & 1877 & Cz & Ln & -- & -- \\
D6NM\footnote{\scriptsize Fine Layered Anorthosite-norite.} & -93 & -25.21080 & 30.09820 & 1868 & Cz & Ln & Clv. 233/16 & \ChFivePhoto{64} \\
MZ5\footnote{\scriptsize No field notes. Assigned Cz because it is stratigraphically below my modeled TU.} & -87 & -25.24180 & 30.07439 & 1829 & Cz & UD & -- & -- \\
D6SN\footnote{\scriptsize Magnetic chromite seam w/ oikocrystic pyroxenite and anorthosite. abundant blocks. chromite hosted with pure anorthosite and in massive layer with oikocrystic pyroxene. Top layers are magnetic (magnetite-bearing?)} & -86 & -25.21040 & 30.09815 & 1863 & UG2 & UG2 & -- & \ChFivePhoto{65}, \ChFivePhoto{66} \\
CZ2 & -83 & -25.21040 & 30.09804 & 1865 & Cz & Ln & -- & -- \\
CZ1 & -72 & -25.20930 & 30.09835 & 1862 & Cz & Ln & -- & -- \\
D4ND\footnote{\scriptsize Coarse Gabbronorite. Indirect measurement, bedding planes in a marsh.} & -72 & -25.21690 & 30.09254 & 1847 & Cz & Pg & Clv. 220/40 & -- \\
CZ & -58 & -25.20720 & 30.09933 & 1884 & Cz & Ln & -- & -- \\
IS & -52 & -25.24350 & 30.07152 & 1840 & IS & IS & -- & -- \\
MZ2\footnote{\scriptsize Assigned Mz because it is undeformed / overlying Kliprivier synform.} & -51 & -25.21740 & 30.09118 & 1847 & Cz & Pg & -- & -- \\
CZ5 & -48 & -25.20630 & 30.09959 & 1898 & Cz & Ln & -- & -- \\
D7SK\footnote{\scriptsize Mottled Anorthosite. Non-magnetic. Interpreted as proximal Merensky Footwall based distance from UG2.} & -33 & -25.20600 & 30.09920 & 1919 & TU & Ma & -- & \ChFivePhoto{81} \\
D11SB\footnote{\scriptsize Porphyritic Gabbronorite. Rotten. Initially confused for mottled anorthosite.} & 162 & -25.22810 & 30.07313 & 1803 & Mz & Pg & -- & -- \\
D11SC\footnote{\scriptsize Mottled Anorthosite. Lower Mottled anorthosite; outcrop source of James' favorite boulder.} & 207 & -25.22700 & 30.07188 & 1806 & LMA & Ma & -- & -- \\
D11NH\footnote{\scriptsize Large cliffside mottled anorthosite outcrop. Modal Layering.} & 333 & -25.19250 & 30.09242 & 1917 & MMA & Ma & M.L. 130/15 & \ChFivePhoto{115} \\
D5NF & 362 & -25.22690 & 30.06477 & 1803 & Mz & UD & Clv. 240/25 & -- \\
D10SJ\footnote{\scriptsize Mottled anorthosite with blebby/blocky recrystalized pyroxenite. Big Outcrop, undulating contacts of MA, often dipping east. Sometimes appear discordant, though others appear concordant: Can be explained by thermal erosion, deposition of surface with meter-scale hummocks and swales.} & 373 & -25.22700 & 30.06419 & 1808 & MMA & MaBap & Clv. 220/15 & \ChFivePhoto{103}--\ChFivePhoto{110} \\
MZ1 & 377 & -25.22870 & 30.06271 & 1797 & Mz & Pg & -- & -- \\
MA & 378 & -25.22920 & 30.06219 & 1767 & MMA & Ma & -- & -- \\
DS & 379 & -25.19100 & 30.09140 & 1903 & Dl & Dl & -- & -- \\
D8SI\footnote{\scriptsize Mottled Anorthosite. Interpreted as LMA. Magnetic mottles. Includes fine-grained gabbronorite with large anhedral CPX phenocrysts ($<$5 cm)} & 398 & -25.18210 & 30.09749 & 1981 & UMA & MaFgog & Clv. 220/20 & \ChFivePhoto{90} \\
D8SH\footnote{\scriptsize Dolerite Sill. Strongly magnetic, with well-developed fine blocky cleavage from stress field. Inherited phenocrystic blocky 1 cm plagioclase with strong ocillitory zoning. Likely later intrusive, not necesarily concordant.} & 413 & -25.18210 & 30.09678 & 1984 & Dl & Dl & Cont. 185/25 & \ChFivePhoto{89} \\
D11NF\footnote{\scriptsize Layered/porhpyritic gabbronorite. Fine modal layering, near "peace" rock.} & 415 & -25.19710 & 30.08507 & 1896 & Mz & Lg & M.L. 130/15 & \ChFivePhoto{113} \\
MA1 & 423 & -25.22600 & 30.06264 & 1805 & MMA & Ma & -- & -- \\
D10SI\footnote{\scriptsize Mottled anorthosite with blebby/blocky recrystalized pyroxenite. MMA.} & 434 & -25.22550 & 30.06250 & 1801 & MMA & MaBap & -- & -- \\
D8NJ & 446 & -25.18260 & 30.09486 & 1979 & Mz & Hg & Clv. 230/15 & -- \\
D11SG\footnote{\scriptsize Mottled Anorthosite. Very large magnetic mottles ($<$10 cm) with large ($<$ 0.5 cm) magnetite.} & 451 & -25.19410 & 30.08572 & 1911 & UMA & Ma & -- & \ChFivePhoto{114}, \ChFivePhoto{117} \\
RESB1ND & 454 & -25.18230 & 30.09475 & 1983 & Mz & Hg & -- & -- \\
D5SG\footnote{\scriptsize Mottled Anorthosite. Float, outcrop in sight. Didn't want to trespass, Also found green fine grained rock with large anhedral CPX.} & 509 & -25.21960 & 30.06343 & 1789 & UMA & MaPxFgog & -- & -- \\
D3SG (D4NC)\footnote{\scriptsize Mottled Anorthosite over strongly graded pyroxenite-anorthosite over fine grained Gabbronorite with large ($<$5 cm) anhedral CPX. Gave a ride to an african man walking home from work at klipriver farms who spoke no english. Stopped by an outcrop and sampled in the dark. Vantongeren maps as Upper Mottled anorthosite.} & 564 & -25.22740 & 30.05494 & 1780 & UMA & MaPxFgog & -- & \ChFivePhoto{49}, \ChFivePhoto{53}, \ChFivePhoto{54} \\
D11NI\footnote{\scriptsize Finely Layered Gabbronorite. Good cleavage, Mottled Anorthosite Nearby?} & 622 & -25.19160 & 30.07954 & 1865 & Mz & Lg & Clv. 240/15 & -- \\
MODALLAYER & 687 & -25.18070 & 30.08540 & 2054 & Mz & Lg & -- & -- \\
D8SG\footnote{\scriptsize Sheared Spotted Anorthosite. Strongy foliated, suggesting co-crystalization shearing. Biggest mottles: $\sim$3cm x 0.5 cm. Interpreted as MMA by RL et al., 2022.} & 708 & -25.18270 & 30.08279 & 2022 & Mz & Sa & -- & \ChFivePhoto{87} \\
MZD & 739 & -25.23480 & 30.04130 & 1790 & Mz & Hg & -- & -- \\
MODAL LAYER & 743 & -25.19140 & 30.07401 & 1828 & Mz & Lg & -- & -- \\
IS2 & 865 & -25.19950 & 30.06205 & 1781 & IS & IS & -- & -- \\
D11SE\footnote{\scriptsize Finely layered gabbronorite. Mottled anorthosite itself is in float here, but graded rocks with blocks are in place.. Mottles are small ($<$3mm) and non-magnetic.} & 867 & -25.18590 & 30.07260 & 1886 & Mz & Lg & -- & \ChFivePhoto{112} \\
D11ND\footnote{\scriptsize $\sim$20 cm Graded spotted anorthosite with irregular contacts and variably recrystalized pyroxenite blocks/blebs. Proper mottled anorthosite in float; non-magnetic. In field interpreted as Upper/main mottled anorthosite (and related package).} & 893 & -25.18530 & 30.07186 & 1891 & Mz & SaLg & -- & \ChFivePhoto{111} \\
D8SC\footnote{\scriptsize Spotted Anorthosite. $\sim$8cm thick, with prominent igneouss layering and including small amounts of pyroxenite below. Interpreted as UMA. I would not call this a mottled anorthosite, the pyroxenes are  do not immediately appear oikocrystic.} & 1030 & -25.17670 & 30.07222 & 1968 & Mz & Sa & M.L. 140/25 & \ChFivePhoto{86} \\
D8NE & 1035 & -25.17780 & 30.07114 & 1955 & Mz & UD & M.L. 130/10 & -- \\
XENOLITH\footnote{\scriptsize A second xenolith near photo 91.} & 1035 & -25.17860 & 30.07051 & 1946 & Mz & Sa & -- & \ChFivePhoto{91} \\
UMA2 & 1035 & -25.17640 & 30.07222 & 1967 & Mz & Sa & -- & -- \\
D8ND\footnote{\scriptsize Layered Gabbronorite. Centimeter-scale igneous layering with variable color and grain size.} & 1036 & -25.17900 & 30.07014 & 1936 & Mz & Lg & M.L. 110/10 & \ChFivePhoto{85} \\
D8ND\footnote{\scriptsize Homogeneous gabbronorite. Different bedding in nearby unlayered rocks...} & 1036 & -25.17900 & 30.07014 & 1936 & Mz & Hg & Clv. 210/10 & -- \\
D8NA\footnote{\scriptsize Porphyritic gabbronorite. Blebby Xenolith noted and photographed.} & 1244 & -25.17050 & 30.06651 & 1814 & Mz & Hg & -- & \ChFivePhoto{84} \\
D6SE\footnote{\scriptsize Green gabbronorite with amphibole "mottles". Amphibole mottles ($<$5 cm) are strongly imbricated, undulose shine, not perfectly optically continuous.} & 1417 & -25.21910 & 30.02170 & 1771 & Mz & A & -- & -- \\
D6SF\footnote{\scriptsize Ironstone concretion. $<$2 cm Goethite nodules enclosing limonite cores with Hematite-Limonite matrix. Rounded clastic quarts $<$2 mm, in float.} & 1417 & -25.21910 & 30.02170 & 1771 & IS & IS & -- & \ChFivePhoto{59} \\
D10NC & 1468 & -25.20220 & 30.03204 & 1755 & Mz & Hg & Clv. 180/10 & -- \\
MZ31 & 1473 & -25.20990 & 30.02600 & 1756 & Mz & Hg & -- & -- \\
D10NB\footnote{\scriptsize Included a range of measured strikes from 250-315, on fault?} & 1475 & -25.20090 & 30.03275 & 1781 & Mz & Hg & Clv. 200/10 & -- \\
D6NB\footnote{\scriptsize Measured another dip (slump?) indirectly from road: 050:30.} & 1551 & -25.21840 & 30.01618 & 1816 & Mz & Hg & Clv. -74/6 & \ChFivePhoto{58} \\
D10ND & 1561 & -25.19970 & 30.02966 & 1772 & Mz & Hg & Clv. 225/15 & -- \\
MZD1 & 1562 & -25.21950 & 30.01472 & 1780 & Mz & Hg & -- & -- \\
D10NA & 1582 & -25.19910 & 30.02914 & 1767 & Mz & Hg & Clv. 200/10 & -- \\
MZD2 & 1623 & -25.22540 & 30.00751 & 1809 & Mz & Hg & -- & -- \\
MZ4 & 1778 & -25.19310 & 30.02448 & 1733 & Mz & Hg & -- & -- \\
D10NH1\footnote{\scriptsize Ironstone Concretion in place.} & 2050 & -25.18090 & 30.02113 & 1746 & IS & IS & -- & -- \\
D10NE & 2056 & -25.17710 & 30.02374 & 1757 & Mz & Hg & Clv. 130/10 & -- \\
MZ12 & 2057 & -25.17750 & 30.02334 & 1740 & Mz & Hg & -- & -- \\
D10NH\footnote{\scriptsize Ironstone Concretion in float.} & 2136 & -25.17770 & 30.01951 & 1729 & IS & IS & -- & -- \\
MZ21 & 2141 & -25.17190 & 30.02370 & 1746 & Mz & Hg & -- & -- \\
D10NG\footnote{\scriptsize Amphibole 'mottled' gabbronorite in float. Fault?} & 2175 & -25.17710 & 30.01813 & 1726 & Mz & A & -- & -- \\
D10NF & 2182 & -25.17820 & 30.01698 & 1730 & Mz & Hg & Clv. 225/15 & -- \\
D6ND\footnote{\scriptsize Dark and Light Facies.} & 2682 & -25.21330 & 29.96669 & 1526 & Mz & Hg & Clv. 160/13 & -- \\
D6NC\footnote{\scriptsize Dark Facies.} & 3266 & -25.19900 & 29.95029 & 1476 & Mz & Hg & Clv. 202/12 & -- \\
\end{longtable}
\end{scriptsize}
\setcounter{footnote}{\ChFiveSavedFootnote}
\endgroup
}
